\begin{claim}{C.44}{Proved}{Exact downstairs forms of a nontrivial boundary quotient}
For the actual quotient $\Phi_z:O\to B$ of C.42,
the open planar orbit $O$ is a plane or an annulus.
Choose a coordinate $B=\R$ when $B$ is a line,
or a coordinate $B=\R/\Z$ when it is a circle,
and use the corresponding target covering in
the definition of $F$. There are source
homeomorphisms giving one of the following
normal forms for the actual map:
\[
\begin{array}{c|c|c|c}
 O & B & \Phi_z & \text{whole fiber}\\ \hline
 \R\times\R & \R & (u,t)\mapsto t & \R\\
 \R\times\R & \R/\Z & (u,t)\mapsto t+\Z
       & \coprod_{\Z}\R\\
 (\R/\Z)\times\R & \R & ([u],t)\mapsto t
       & \R/\Z\\
 (\R/\Z)\times\R & \R/\Z & ([u],t)\mapsto t+\Z
       & \coprod_{\Z}(\R/\Z)\\
 \R\times(\R/m\Z) & \R/\Z & (u,[t]_m)\mapsto[t]_1
       & \coprod_{j=0}^{m-1}\R,\quad m\geq1 .
\end{array}                                             \tag{C.44.1}
\]
In the last row the actual induced image
$\Phi_{z*}(\pi_1O)$ is $m\Z\subset\pi_1B=\Z$;
in all other rows that image is zero.
Every downstairs fiber component is a
properly embedded line or a circle, and
the source covering restricts onto each
such component as the line covering
specified by these formulas.
\end{claim}

\begin{proof}
Let $\Gamma$ be the deck group of
$p:\widetilde O\to O$.
Use C.43 to identify the covering plane
with $\R_u\times\R_t$ so that $F(u,t)=t$.
The equation $qF=\Phi_zp$ implies, for
each $\gamma\in\Gamma$,
\[
 F\gamma=F+\eta(\gamma),                                \tag{C.44.2}
\]
where $\eta(\gamma)=0$ for a line target
and $\eta(\gamma)\in\Z$ for a circle
target. For the circle, the difference
is a continuous integer-valued function
on the connected covering plane and
is therefore constant. Composition
shows that $\eta:\Gamma\to\Z$ is a
homomorphism. Its value is exactly the
winding number of the image under
$\Phi_z$ of the source loop whose lift
ends at $\gamma$ times its initial
point. Thus its image is the actual
induced fundamental-group image.

Write $K=\ker\eta$. Every deck map
preserves the orientation pulled back
from the planar surface $O$. It has
the form
\[
 \gamma(u,t)=
       (h_{\gamma,t}(u),\,t+\eta(\gamma)),                \tag{C.44.3}
\]
and each $h_{\gamma,t}$ is increasing:
the base coordinate is increasing,
and reversing the fiber direction
would reverse the surface orientation.
The functions in (C.44.3) and their
fiberwise inverses are jointly
continuous, by the deck map and its
inverse.

\emph{The exact kernel on a real fiber.}
The action of $K$ on any fiber
$\R\times\{t\}$ is free and proper.
Freeness follows from deck freeness.
For properness, compact subsets of
this fiber are compact in the
covering plane, and the deck action
has finite transporters of compact
pairs. The restriction is faithful,
since a deck transformation fixing
one point is the identity.
An increasing free proper group on
a line is trivial or infinite cyclic.
Here is the direct argument.
Its orbit of a point is closed and
discrete. In the nontrivial case
choose the element taking that point
to its nearest orbit point on the
right. Its iterates tend to the
two ends, because a finite limiting
point would be fixed. Every other
orbit point is one of these iterates:
an intervening point could be
translated back between the original
point and its nearest right neighbor.
Freeness then identifies every
element with the corresponding power.
Consequently
\[
                          K=1\quad\hbox{or}\quad K\cong\Z.
                                                               \tag{C.44.4}
\]

Suppose $K=\langle k\rangle\cong\Z$,
choosing the generator to move one
fiber point to the right.
Since the displacement
$h_{k,t}(u)-u$ is continuous and
nowhere zero on the connected plane,
it is positive everywhere.
Put $k_t=h_{k,t}$.
For $v=n+r$, where $n\in\Z$ and
$0\leq r\leq1$, define
\[
 H(v,t)=\bigl(k_t^n(r\,k_t(0)),\,t\bigr).
                                                               \tag{C.44.5}
\]
At integral $v$ the two formulas
agree. The intervals with endpoints
$k_t^n(0)$ exhaust the line in both
directions, so $H$ is a continuous
bijection of the plane. Invariance
of domain makes it a homeomorphism.
It preserves $F$ and satisfies
\[
                        H(v+1,t)=kH(v,t).
                                                               \tag{C.44.6}
\]
Thus the actual intermediate quotient
is
\[
 \widetilde O/K\cong(\R/\Z)\times\R,
 \qquad F_K([v],t)=t.                                  \tag{C.44.7}
\]

\emph{The compactness obstruction retaining the quotient map.}
The case $K\cong\Z$ cannot coexist
with a nonzero image of $\eta$.
Indeed that image would be $m\Z$
for some $m\geq1$. The quotient
$\Gamma/K$ is then infinite cyclic,
and its chosen generator acts on
(C.44.7) by a circle homeomorphism
in the first coordinate and by
$t\mapsto t+m$ in the second.
The compact strip
$(\R/\Z)\times[0,m]$ maps onto the
full quotient by $\Gamma/K$:
each point can be translated into
that strip using its second coordinate.
It follows that
\[
     O=\widetilde O/\Gamma
       =(\widetilde O/K)/(\Gamma/K)
\]
is compact, contrary to its being
a nonempty open subset of the plane.
Therefore $K\cong\Z$ forces
$\eta=0$ and $\Gamma=K$.
Equations (C.44.5)--(C.44.7) give
the third and fourth rows of
(C.44.1), according to the target.

If $K=1$ and $\eta=0$, then
$\Gamma=1$, so C.43 gives the
first two rows.

\emph{The remaining nonzero covering degree.}
It remains that $K=1$ and
$\eta(\Gamma)=m\Z$, with $m\geq1$.
Then $\Gamma$ itself is infinite
cyclic. Choose its generator
$\gamma$ with $\eta(\gamma)=m$,
and write
\[
                   \gamma(u,t)=(h_t(u),t+m).
                                                               \tag{C.44.8}
\]
We give a coordinate change
preserving the real target
coordinate and removing $h_t$.
For $0\leq t\leq m$ set
\[
 \rho_t(u)=
          \left(1-\frac{t}{m}\right)u
                   +\frac{t}{m}h_0(u).
                                                               \tag{C.44.9}
\]
Each $\rho_t$ is an increasing
homeomorphism of the line, with
$\rho_0=\id$ and $\rho_m=h_0$.
Extend to all real $t$ in both
directions by the exact recursion
\[
                     \rho_{t+m}=h_t\circ\rho_t.
                                                               \tag{C.44.10}
\]
The endpoint identity makes the
extension jointly continuous at
every multiple of $m$.
The map
$H_1(u,t)=(\rho_t(u),t)$ is a
continuous bijection of the plane
and hence a homeomorphism.
Equations (C.44.8)--(C.44.10) give
\[
              H_1^{-1}\gamma H_1(u,t)=(u,t+m).
                                                               \tag{C.44.11}
\]
The quotient is therefore
$\R\times(\R/m\Z)$, and the
original map, still satisfying
$qF=\Phi_zp$, is precisely
$(u,[t]_m)\mapsto[t]_1$.
This is the final row of (C.44.1).

All changes of source coordinates
preserved $F$; the target was
changed only by its initially
specified coordinate.
The whole-fiber descriptions now
follow directly. For a circle
target, levels $F=t_0+j$ are
indexed by $j\in\Z$. The deck
group identifies their labels
by $\eta(\Gamma)$ and acts
within each level by $K$.
Thus the component labels are
$\Z/\eta(\Gamma)$ and each
component is $\R/K$ in the
explicit coordinates above.
For a line target there is
only one real level and the
component is again $\R/K$.
These formulas also prove
surjectivity of the restricted
source covering onto each
downstairs component and give
its deck group. In particular,
they establish the component
assertion without assuming
component-surjectivity for an
arbitrary restricted covering.
\end{proof}

\paragraph{Scope of the reduction.}
C.44 classifies each of the actual
quotient maps separately, including
its whole fibers and covering
labels. Several boundary points
can impose several such maps on
the same open orbit. No common
product coordinates or homotopy
preserving all those maps have
been constructed here. That
simultaneous compatibility and
frontier extension remain part
of C.12.
